\section{长时序建模：CNN + RNN 到 self-attention}

\subsection{为什么短 clip 不够}

短 clip 能捕获局部动作，但很多任务需要更长时间跨度上的语义。例如，一个人先站、后跑、再停下，或者一个动作从准备、执行到完成，往往都需要跨越更长的时间窗口。此时只看几秒钟就不够了。

\begin{figure}[H]
\centering
\includegraphics[width=0.95\textwidth]{slides-images/slide-049.jpg}
\caption{短 clip 只能覆盖局部运动：如果动作跨越更长时间，就需要长时序建模\protect\footnotemark}
\end{figure}
\footnotetext{Slides 第50页。}

\subsection{CNN + RNN / LSTM}

一个自然的想法是：先用 CNN 对每个短片段或每一帧提取特征，再把这些特征送进 RNN 或 LSTM 聚合时间信息。这样做的好处是把问题拆成两个层次：

\begin{itemize}
    \item CNN 负责局部空间或局部时空特征；
    \item RNN 负责更长范围的时间依赖。
\end{itemize}

\begin{figure}[H]
\centering
\includegraphics[width=0.95\textwidth]{slides-images/slide-050.jpg}
\caption{长时序建模的第一步：先用 CNN 提取每个时间片的局部特征，再交给序列模型处理更长的依赖\protect\footnotemark}
\end{figure}
\footnotetext{Slides 第51页。}

\begin{figure}[H]
\centering
\includegraphics[width=0.95\textwidth]{slides-images/slide-051.jpg}
\caption{如果用 LSTM 处理这些局部特征，就能把短时序表征扩展到更长的视频级表示\protect\footnotemark}
\end{figure}
\footnotetext{Slides 第52页。}

\begin{knowledgebox}{CNN + RNN 的核心思想}
CNN 看到的是局部窗口，RNN 看到的是时间序列。把两者串起来，就能兼顾“局部动作模式”和“全局视频语义”。
\end{knowledgebox}

\subsection{Recurrent convolutional network}

课程还提到一种更细粒度的组合方式：把 RNN 的矩阵乘法替换成卷积，得到 recurrent convolutional network。它的思想是：

\begin{itemize}
    \item 在空间维度上保留卷积的局部归纳偏置；
    \item 在时间维度上递归地传递状态；
    \item 通过时空联合建模来处理视频序列。
\end{itemize}

这种设计在概念上很自然，但受限于 RNN 本身的顺序性，训练和推理都不算高效。

\begin{importantbox}{为什么长序列会逼出 attention}
RNN 可以做长时序建模，但它难并行、难优化、长距离依赖也容易变弱。attention 的价值，就是把“远处帧之间可以直接连接”这件事做成显式计算。
\end{importantbox}

\subsection{self-attention 和 non-local block}

当视频足够长时，局部卷积和循环结构仍然会受限。一个更直接的思路是引入 \textbf{self-attention} 或 \textbf{non-local block}：让任意两个时空位置都能交互，而不是只看邻近窗口。

如果输入向量为 $x$，那经典 self-attention 可以写成：

\[
q = xW_q,\quad k = xW_k,\quad v = xW_v,
\]
\[
e_{ij} = \frac{q_i^\top k_j}{\sqrt{d}}, \qquad a_{ij} = \mathrm{softmax}(e_{ij}), \qquad y_i = \sum_j a_{ij} v_j.
\]

\begin{itemize}
    \item Query / Key / Value 的逻辑和 Transformer 类似；
    \item 它可以在 3D 特征图上做全局匹配；
    \item 这样能捕获更长范围的时空依赖。
\end{itemize}

\begin{knowledgebox}{self-attention 在视频里有什么用}
\begin{itemize}
    \item 它让远距离帧之间直接交互。
    \item 它适合捕获长程依赖和非局部关系。
    \item 它为后续 Transformer 化的视频模型铺路。
\end{itemize}
\end{knowledgebox}

\begin{importantbox}{长时序建模的共同目标}
不管是 RNN、Recurrent CNN、Non-local block 还是 Transformer，目标都一样：把“很远的帧之间可能有关联”这件事显式编码进模型。
\end{importantbox}

